% Pro Cn. Plancio --- headnote
% pro-plancio, 54 BC, Rome

\textit{For Cn. Plancius, delivered at Rome in 54 BC.
Plancius, the son of a leading Roman knight and
tax-farmer from Atina in the Volscian country, had just
been elected curule aedile and was now prosecuted for
ambitus --- electoral corruption --- under the lex
Licinia de sodaliciis, the law of 55 BC against the
organized electoral clubs. The prosecutor was M.
Iuventius Laterensis, a young noble who had stood for
the same aedileship and lost; his subscriptor was L.
Cassius Longinus. The case was, for Cicero, a debt of
honour. As quaestor in Macedonia in 58 BC, Plancius had
taken in the exiled Cicero at Thessalonica and sheltered
him through the worst months of his banishment; the
speech is suffused with that gratitude, and Cicero says
plainly that he must plead not only for Plancius but
for himself, since the prosecution had spent more words
on him than on the defendant.}

\textit{The law on sodalicia gave the prosecutor an
unusual weapon: the right to name (\textit{edere}) the
tribes from which the jury would be drawn --- the
\textit{editicii iudices} --- on the theory that the
bribery worked tribe by tribe. Cicero's legal answer,
in the central sections, is that Laterensis has charged
the wrong offence: there is no trace of the organized
tribe-buying the law was written to punish, only the
ordinary friendships and home-town loyalties by which an
election is honestly won. Around this turns the speech's
most famous matter --- the long reflection on why one
man is returned at the polls and another is not.
Electoral defeat, Cicero insists, is no verdict on a
man's worth: the steps of office stand open equally to
the highest and the lowest, but the steps of glory are
unequal (\textsection 60), and Laterensis, who lost,
has no standing to read his own loss as Plancius's
crime. The comparison of the two men runs through their
families and their towns --- the modest municipal
dignity of Atina set beside the older fame of Tusculum
--- and through the contrasting tempers that carried the
one candidate and sank the other.}

\textit{The closing movement returns to Cicero's exile.
Plancius's loyalty at Thessalonica becomes the moral
centre of the defence, and the peroration is among
Cicero's most personal: the recollection of his
banishment, the plea that the man who guarded his life
should not be destroyed for the very act of guarding it,
and the appeal to the jurors' own memory of what they
owed and what was owed to them. Plancius was acquitted.}
